\section{Discussion}\label{sec:discussion}
\subsection{What is now answered}
Our main question was whether MICA's reported advantage survives when each method must
choose its own clustering without reference labels. On the evidence here the honest answer
is that the question, as we posed it, does not have a single answer, and the reason is
informative rather than evasive.

Under one published rule read exactly as its authors specify it, MICA is ahead. Under the
same scores read with a different tie-break, it is behind. Nothing about the data or either
clustering method changes between those two readings. What changes is how the rule resolves
settings it cannot tell apart, and for MICA it cannot tell apart a large fraction of them.
A benchmark that reports a single label-free rule can therefore obtain either ranking, and
would in both cases appear to be making a statement about the methods.

This is a different claim from the one we set out to test, and weaker in one sense and
stronger in another. Weaker, because we cannot say which method a user should prefer.
Stronger, because it identifies a failure mode that applies to any comparison of this kind,
not only to ours.

\subsection{Why it happens}
The mechanism we can support is that MICA's partitions are too reproducible for
reproducibility to be a useful criterion. MICA's small-data pipeline is deterministic given
the distance matrix, and its distance is unusually stable under subsampling because it is
computed per cell pair rather than from any population quantity. Splitting a true cell type
into several clusters is then just as repeatable across subsamples as recovering it, so the
score barely distinguishes the two. The baseline rebuilds its principal components and its
neighbour graph on every subsample, which injects exactly the variability the rule is
designed to detect.

We tested this rather than assuming it. Two predictions follow and both hold. First, the
span of cluster counts among settings that tie with the best should be wider for MICA, and
it is. Second, where that ambiguity is widest the rule's preference for fine partitions
should produce the worst overclustering, and the association runs in the predicted
direction. We note one competing explanation we could not rule out: all but one of our
datasets has fewer than 3,000 cells, while chooseR was demonstrated near 11,000, so some of
this may be a small-sample effect rather than a property of mutual information clustering.
Separating the two needs a dataset where MICA's graph mode and small-data mode can both be
run, which is Section~\ref{sec:conclusion}'s first piece of remaining work.

\subsection{What it changes for someone using these tools}
Three practical consequences follow, each bounded by what we tested.

A stability rule should not be used on its own to choose the number of clusters for MICA on
datasets of this size. It will return a defensible-looking score for a partition with
several times too many clusters. Reporting the chosen cluster count alongside the score
makes this visible immediately, which is why we show it.

The gap between a method's benchmark score and what a user obtains is not a constant that
can be quoted once. It depends on the selection rule, and on these datasets the spread
across rules is larger than the gap between the two methods under any single rule. A
benchmark that fixes settings with the labels is not measuring a quantity that transfers to
practice.

Finally, MICA's small-data mode returns four clusters when run unmodified, whatever the
dataset. Nine of the ten benchmark datasets fall in that mode. Any comparison that does not
say how the number of clusters was obtained is, for those datasets, reporting a number that
came from outside the method.

\subsection{Limitations}
The comparison covers one clustering method against one baseline, under rules we chose,
on ten datasets from a single published benchmark. Nine of the ten are smaller than the
data chooseR was demonstrated on, and the tenth is the only one in MICA's graph mode, so
mode and size are confounded in our design and we cannot separate them. Three of the
datasets carry labels derived by clustering the same data, which flatters methods that
resemble the original pipeline. We report an observed range over three seeds rather than a
confidence interval, because three runs do not support one. The bootstrap inside the rule is
a percentile bootstrap where chooseR specifies BCa, which affects the threshold and in
principle the choice, although we saw no case where it did. Our implementation of the rule
reproduces chooseR's procedure from its description rather than running the authors' R code,
so a discrepancy in our reading of it would propagate to every result in
Table~\ref{tab:rules}.
